\chapter{Implementation Guide}
\label{app:implementation_guide}
Chapter~\ref{ch:geometric_predicates} defines the exact-sign predicates that
underpin the algorithms in this book; this appendix is the practical
counterpart. It explains how to choose a numerical representation, how to
implement the predicates so their signs are trustworthy, how to lay out
memory for speed, how to test code that must fail rarely and crash never,
and how to debug the failures that slip through. The advice is opinionated:
for most applications the winning configuration is double precision for
storage, exact-sign predicates protected by floating-point filters, and an
extensive battery of degenerate and near-degenerate tests.
\section{Designing a Geometry Kernel}
\label{sec:app_impl_kernel}
A \index{kernel!geometry}geometry kernel is the small set of types and
operations every algorithm in the library is built on. Keep it small: a
kernel with a dozen ``point-like'' types will contain at least one pair no
one remembers is different, and the bugs will live in the conversions. Three
structural decisions matter more than any individual function.
\subsection{Predicate-First Design}
\label{sec:appimpl_predicate_first}
Every algorithm in this book reduces to a handful of sign tests. Design the
kernel around those tests and let constructed objects (points, segments,
triangles) be plain data:

\begin{itemize}
  \item Predicates return the sign $-,0,+$ (an \texttt{int} or enum), never
        a distance or coordinate. Returning a determinant so callers
        ``compare it against zero themselves'' is how epsilon thresholds
        sneak in and rob the library of consistency.
  \item Predicates take indices or pointers into coordinate arrays, not
        freshly constructed point objects, so the hot path pays no
        allocation or virtual dispatch.
  \item Provide the canonical family $\operatorname{orient2d}$,
        $\operatorname{orient3d}$, and $\operatorname{incircle}$, matching
        Sections~\ref{sec:orientation_predicate},
        \ref{sec:orientation3d_predicate}, and \ref{sec:incircle_predicate}.
        Express everything else (between, on-segment, left-turn) through
        these three so there is one sign definition to audit.
  \item When an algorithm needs a real number --- an intersection point, a
        Voronoi vertex --- construct it in a separate, clearly named layer.
        This is the seam at which exactness is lost, and it should be
        visible.
\end{itemize}
\subsection{Memory Layout and Cache Locality}
\label{sec:appimpl_layout}
For large inputs, geometric code is usually memory-bound: a predicate costs
a few floating-point operations, and a modern core executes hundreds in the
time it takes to miss the cache. Layout therefore frequently dominates
wall-clock performance~\cite{ericson2004}.

\index{array of structures}\textbf{Array-of-structures (AoS)} stores each
point as one struct in a contiguous array; use it when the algorithm reads
several fields of the \emph{same} point, as in traversals and walks.
\index{structure of arrays}\textbf{Structure-of-arrays (SoA)} keeps one
array per coordinate (\texttt{xs[]}, \texttt{ys[]}, \texttt{zs[]}); use it
when the code sweeps many points doing the same operation on one field, as
in bounding-box rejection. With an inline coordinate accessor both are
easy to switch, so treat layout as a late, measured decision.

Two rules apply to both. First, \index{cache locality}iterate in the order
data was written: random access into a large array of arbitrary points is
the most cache-hostile pattern in the field, and preprocessing (sorting,
blocking, space-filling curves) often turns a cold loop into a warm one.
Second, pack hot records to a cache line; a half-edge record padded past 64
bytes wastes bandwidth on every traversal.
\subsection{Half-Edge Layout: Index-Based versus Pointer-Based}
\label{sec:appimpl_halfedge_layout}
Section~\ref{sec:halfedge_comparison} compares representations at the
interface level; the implementation choice between indices and pointers is
driven by how the mesh grows~\cite{ericson2004}:

\begin{itemize}
  \item \index{half-edge!index-based}\textbf{Index-based} half-edges store
        \texttt{next}, \texttt{prev}, \texttt{twin} as integers into one
        array of records. Copies are cheap, records are relocatable, and
        growth never invalidates references. The best default for
        incremental construction, streaming, and editing.
  \item \textbf{Pointer-based} half-edges traverse marginally faster, but
        reallocation invalidates every pointer, so the mesh must be fully
        built before it is exposed.
  \item A common hybrid owns by index and caches
        \texttt{next}/\texttt{twin} pointers rebuilt after growth; the
        pointers pay for themselves only when traversal is measurably hot.
  \item \index{64-bit index}Use 64-bit index types once the mesh can exceed
        roughly two billion edges; each reference costs one extra word, so
        keep 32-bit indices for meshes known to be small.
\end{itemize}

The storage analysis of Section~\ref{sec:halfedge_representation} and the
incidence identities of Theorem~\ref{thm:hed_incidence} are the invariants
a well-tested layout must preserve regardless of encoding.
\subsection{Integration Patterns and Thread Safety}
\label{sec:appimpl_integration}
Different callers impose different requirements on the same kernel:
\index{batch processing}\textbf{Batch} (offline) callers know all input up
front: preallocate capacities, run in sorted order, never pay for dynamic
growth --- the fastest configuration and the one to benchmark.
\textbf{Streaming} callers supply elements continuously: prefer
index-based references so growth never invalidates handles, and bound
reallocation cost with geometric growth factors. \textbf{Interactive}
callers edit between queries: keep an undo log of primitive operations
(split, flip, insert) rather than re-snapshot meshes, and run invariant
checks after each edit in debug builds.

\textbf{Thread safety.} The common pattern is to build the structure on one
thread, \index{immutable structure}then publish it immutable and let many
threads read concurrently --- immutable structures are automatically
thread-safe. Predicates that read only their arguments and write to
per-thread scratch space are thread-safe; predicates that cache in global
state are not. For incremental construction, serialize mutation or
partition the domain so each thread owns a disjoint region, and freeze the
kernel's numeric backend before the first parallel query.
\section{Choosing Numeric Types}
\label{sec:app_impl_numeric_types}
The first decision in any geometric library is which number type the kernel
is built on, and it is the most expensive one to change later, so it should
be made consciously rather than by default.
\subsection{The Trade-Off at a Glance}
\label{sec:appimpl_representations}
\begin{center}
\begin{tabular}{p{3.2cm}p{2.2cm}p{2.2cm}p{3.0cm}p{4.0cm}}
\hline
\textbf{Type} & \textbf{Exact?} & \textbf{Speed} & \textbf{Input} & \textbf{Choose when} \\
\hline
\index{double precision}double & no & fastest & floats &
measurements; tolerances acceptable \\
single precision & no & very fast & floats & memory-bound 3D meshes \\
64-bit integer & yes & fast & integer coords & indices, grids, integer
data \\
\index{rational arithmetic}rational & yes & moderate & integer/ratio &
construction-heavy pipelines \\
arbitrary precision & yes & slow & any & bit growth unavoidable; fallback \\
\index{expansion arithmetic}expansions & yes & fast* & doubles &
filtered predicates (Section~\ref{sec:appimpl_adaptive}) \\
\hline
\end{tabular}
\end{center}

*Fast in the common case; the exact fallback is slow. The starred entry is
where robust libraries converge: ordinary doubles for storage, exact
arithmetic only where a sign must be certified.
\subsection{Double Precision First}
\label{sec:appimpl_double}
For measurements, sensor data, or coordinates that carry their own
uncertainty, \index{double precision}double precision is the right default:
it is hardware-supported, and any value more precise than the data is
wasted. The discipline is that \emph{constructed} values are approximate:
an intersection point computed in doubles has error on the order of the
input's rounding error times a modest constant~\cite{higham2002,goldberg1991}.
Use doubles when the algorithm must answer ``is this configuration
plausible''; escalate when it must answer ``is this combinatorial relation
true, exactly.''
\subsection{Exact Integer and Rational Kernels}
\label{sec:appimpl_integer}
When input coordinates are integers or rationals, exactness is cheap and
should be used. The key fact (Section~\ref{sec:exact_arithmetic}) is that
predicate bit growth is bounded: for $\tau$-bit coordinates the 2D
orientation determinant needs at most $2\tau+1$ bits and the in-circle
determinant $O(\tau)$ bits, so 64-bit integers cover $\tau \le 31$ and
arbitrary-precision integers cover the rest with predictable cost. Rational
kernels add exact construction of new points
(Section~\ref{sec:rational_arithmetic}); their cost is acceptable for
engineering input, as the Delaunay experiments of Karasick, Lieber, and
Nackman showed~\cite{karasick1991}. Fortune and Van Wyk's expression
compiler pushes this further, generating unrolled evaluation code whose cost
tracks the operand bit-length instead of general library
overhead~\cite{fortune1993}.
\subsection{The Cost of Exactness and Arbitrary-Precision Fallbacks}
\label{sec:appimpl_cost}
\index{arbitrary precision}Arbitrary-precision arithmetic is the correct
fallback when constructed points keep growing, as in chains of intersection
operations: the exact-sign framework of Section~\ref{sec:exact_arithmetic}
shows the bit-length grows only by a constant per predicate, but a planar
subdivision's coordinates grow with the depth of construction. Reserve it
for the exact kernel behind a filter, never for the whole library.

Measured costs are consistent: a naive exact predicate is one to two orders
of magnitude slower than its floating-point twin, and an arbitrary-precision
construction can be far worse. Users rarely see that cost, because an
\index{floating-point filter}exact sign is needed only for nearly degenerate
inputs, which are rare. A filtered predicate runs at roughly one to three
times the speed of the plain floating-point version in the common
case~\cite{shewchuk1997,bronnimann2001}. Spend the design effort on the
filter, not on making the exact kernel fast.
\section{Floating-Point Error}
\label{sec:app_impl_floating_point}
Predicates fail when rounding flips a sign. Understanding the size and the
\emph{structure} of that error is prerequisite to designing filters, and
the most useful numerical knowledge a geometric programmer needs
~\cite{goldberg1991,higham2002}.
\subsection{Machine Epsilon and Relative Error}
\label{sec:appimpl_machine_eps}
IEEE~754 doubles have a 53-bit significand and machine epsilon
$\varepsilon = 2^{-53} \approx 1.11\times 10^{-16}$; each elementary
operation is correctly rounded, so its relative error is at most
$\varepsilon/2$. Error bounds are conventionally $k\varepsilon$ times a sum
of magnitudes of intermediate values. The number of operations matters far
less than the \emph{magnitudes}: a determinant of large coordinates
accumulates absolute error proportional to the largest product, while its
exact value may be tiny.
\subsection{Where Geometric Computations Lose Accuracy}
\label{sec:appimpl_loss_accuracy}
Two effects dominate. \index{catastrophic cancellation}\textbf{Catastrophic
cancellation} occurs when nearly equal values are subtracted: the leading
digits cancel, the result is almost pure roundoff, and the \emph{relative}
error explodes while the absolute error stays small. In the orientation
predicate
\[
  \operatorname{orient2d}(A,B,C) =
  (B_x-A_x)(C_y-A_y) - (B_y-A_y)(C_x-A_x),
\]
cancellation is mild for well-separated points but total for nearly
collinear ones. Second, \textbf{error accumulation} compounds through
intermediate steps: subtracting coordinates then multiplying roughly doubles
the error bits, so a two-term product carries error of order a few
$\varepsilon$ times its magnitude. This is the bound a static filter
certifies against.
\subsection{What Epsilon Can and Cannot Do}
\label{sec:appimpl_epsilon}
An epsilon threshold is a tolerance, not a correctness mechanism. It can
legitimately express an application's notion of ``good enough'' --- snapping
nearby vertices, merging close segments for display --- and it is the only
tool for genuinely approximate input. It cannot provide a \emph{consistent
combinatorial result}: a threshold that absorbs one near-degenerate case
makes an algorithm accept a configuration that a neighboring, equally valid
case rejects, and topology can change under rotation. As the robustness
literature stresses (Section~\ref{sec:robustness_failures}), an algorithm
whose decisions rest on thresholds can loop forever, crash, or emit
inconsistent output that a consistent sign-based algorithm never
would~\cite{kettner2008}.

A practical hygiene checklist:

\begin{itemize}
  \item Never test \texttt{a == b} or \texttt{a <= 0} on a computed value
        that involved subtraction. Compare against a \emph{relative}
        tolerance derived from the data's scale, or better, use an exact
        predicate.
  \item When a tolerance is unavoidable, scale it by the operand magnitude
        ($\varepsilon_{\mathrm{tol}} \approx k\varepsilon \cdot
        \max(|a|,|b|)$) rather than using an absolute constant.
  \item Do not write $\varepsilon$ into predicates that drive topology;
        route topology decisions through the exact kernel.
  \item If two computations must agree on which side of a line a point
        lies, both must use the \emph{same} predicate --- recomputing with a
        differently ordered formula invites disagreement.
\end{itemize}
\section{Exact Predicates}
\label{sec:app_impl_exact_predicates}
This section covers the two filter families and the adaptive machinery
behind them (the theory is in Section~\ref{sec:floating_point_filters}):
robustness comes from making the predicate sign exact, and the cheapest
route makes it exact in the common case while retaining an exact fallback.
\subsection{Only the Sign Matters}
\label{sec:appimpl_sign}
Every determinant-based predicate is a device for producing a sign; the
magnitude is irrelevant to the algorithm, should never be consumed by it,
and should never feed back into a construction (it carries the same roundoff
as any computed float). The exact fallback therefore needs to compute only
enough to fix the sign --- exactly what an \index{adaptive precision}adaptive
algorithm does: each stage is a more accurate estimate, and the first stage
whose sign is certified wins.
\subsection{Static Filters}
\label{sec:appimpl_static}
\index{floating-point filter!static}A \textbf{static filter} bounds the
 roundoff before the computation. For $\operatorname{orient2d}$ the classic
 bound is
\[
  |\mathrm{error}| \le 3\varepsilon\,\bigl(
  |(B_x-A_x)(C_y-A_y)| + |(B_y-A_y)(C_x-A_x)|\bigr),
\]
so the sign is certified whenever the computed determinant exceeds that
bound --- one or two fused operations, then ``if clearly nonzero, trust
it.'' Shewchuk's public-domain predicates ship exactly this way: a few
floating-point checks in the fast path, exact arithmetic only in the
tail~\cite{shewchuk1997}. The bound is pessimistic, but that is the price
of a fixed, one-pass test.
\subsection{Dynamic Filters and Interval Arithmetic}
\label{sec:appimpl_dynamic}
\index{floating-point filter!dynamic}A \textbf{dynamic filter} tightens the
bound as the evaluation proceeds, typically wrapping each operation in an
\index{interval arithmetic}interval so the result carries a running error
interval: if the interval stays clear of zero, the sign is certified with
the true error, not an a priori bound. Dynamic filters resolve a
substantially larger fraction of cases than static filters at modest extra
cost per operation~\cite{bronnimann2001}. The adaptive predicates below are
the stronger version of the same idea: instead of intervals, recompute the
determinant with progressively more accurate error-free arithmetic.
\subsection{Adaptive Precision and Expansion Arithmetic}
\label{sec:appimpl_adaptive}
\index{expansion arithmetic}Expansion arithmetic represents an exact value
as the sum of non-overlapping floating-point components. Its primitives are
the error-free transformations \index{two-sum}\textbf{two-sum} and
\index{two-product}\textbf{two-product}: given IEEE~754 inputs each returns
the rounded result \emph{and} the exact rounding error, so an exact sum of
arbitrary precision is built from ordinary floating-point operations
(``What Every Computer Scientist...'' explains why the rounding is
recoverable~\cite{goldberg1991,priest1991}). Two-sum, for example, is four
operations:

\begin{verbatim}
s = a + b
bp = s - b
ap = s - bp
err = (a - ap) + (b - bp)
\end{verbatim}

Here \texttt{err} is the exact error of the addition. Adaptive predicates
rebuild a determinant as an expansion of length $3, 5, 7, \dots$ components,
stopping at the first stage whose sign is certified; the exact stage on
double inputs never needs more than a few dozen components, and the common
case never leaves the first stage~\cite{shewchuk1997}. The result is the
guarantee of exact arithmetic at a fraction of its cost, which is why the
technique is the de facto standard. The same primitives compose into full
exact arithmetic when a constructed value, not just a sign, must be exact
(Section~\ref{sec:exact_arithmetic}).
\subsection{Handling Degeneracies}
\label{sec:appimpl_degeneracies}
Degenerate configurations --- \index{collinear}collinear points,
\index{cocircular}cocircular points, coincident points --- are ordinary
inputs a correct library must answer, not edge cases to paper over. Three
strategies exist, and they can be combined:

\begin{itemize}
  \item \index{general position}\textbf{General position as a contract.}
        Many algorithms are proved only for general position. Document the
        assumption, \emph{verify} it once at input time (scan for
        duplicates, collinear triples, cocircular quadruples), then run the
        fast path. Rejecting or snapping input is an application decision
        and must not silently change an algorithm's answer.
  \item \index{simulation of simplicity}\textbf{Symbolic perturbation.}
        Simulation of Simplicity perturbs the input infinitesimally --- in
        the predicates only, never in the data --- so tie-breaks resolve
        consistently, in the manner of
        Section~\ref{sec:simulation_of_simplicity}~\cite{edelsbrunner1990}.
        It is elegant and cheap, but the tie-break order is symbolic, not
        geometric: different perturbations yield different but internally
        consistent answers.
  \item \textbf{Explicit handling.} Deduplicate identical points, treat zero
        orientation as an event (a point on a segment, a shared vertex),
        and give the algorithm well-defined tie-break rules
        (Section~\ref{sec:appimpl_pitfalls}). The most work, and the only
        option when the degenerate answer carries geometric meaning, e.g.,
        a mesh that must stay manifold across a shared edge.
\end{itemize}
\section{Testing Geometric Software}
\label{sec:app_impl_testing}
Geometric code fails in distinctive ways: wrong answers are rare, but when
they occur they are crashes, hangs, or subtly inconsistent structures, and
one bad predicate decision corrupts everything downstream. The strategy
below is built around that reality.
\subsection{Randomized and Property-Based Testing}
\label{sec:appimpl_randomized}
\index{randomized testing}Random inputs with a fixed seed find bugs
hand-written examples never do, and the fixed seed makes failures
reproducible. \index{property-based testing}Property-based testing pushes
this further: the generator enumerates cases and the test asserts a
\emph{property} that must hold for all of them. Properties for geometry
include: the convex hull of any point set is convex and contains every input
point; a triangulation satisfies Euler's formula $V-E+F=2$, and the
half-edge structure satisfies $\operatorname{twin}(\operatorname{twin}(e))
= e$ and the identities of Theorem~\ref{thm:hed_incidence}; flipping a
Delaunay edge never destroys the triangulation property; and a point-location
query agrees with a brute-force scan.

Generators must hit \emph{structure}: uniform points in a square, points on
a grid, points on a circle, degenerate families, and coordinates spanning
many orders of magnitude. Uniform-random inputs almost never produce a
collinear triple, so a suite that draws only uniform points is silently
testing nothing about degeneracies.
\subsection{Brute-Force Oracles for Small Inputs}
\label{sec:appimpl_bruteforce}
The most powerful oracle is a second, obviously-correct implementation
written only for tiny inputs. Compare your fast algorithm against it on
every random case of size $n \le 30$: check the Delaunay in-circle
predicate against an $O(n^4)$ scan over all quadruples (or a reference
library such as the one surveyed in Section~\ref{sec:app_soft_cgal});
verify every convex-hull edge is a supporting line; compare the reported
segment intersections against an $O(n^2)$ test using an exact predicate;
and compare point-in-polygon against an \index{oracle}oracle built from the
winding-number definition (Section~\ref{sec:point_in_polygon_winding}). The
oracle is short and slow, so it is cheap to audit by hand; keep it
permanently as a ``reference mode'' rather than deleting it once the
algorithm ``passes.''
\subsection{Invariant Checking}
\label{sec:appimpl_invariant}
\index{invariant}Invariants are the geometry library's analog of
assertions. Run them on every mutation in debug builds and on random
snapshots in release builds: \index{Euler's formula}Euler/incidence
($V-E+F=2$ for a connected planar subdivision, and every half-edge's
twin/next chains close); manifoldness (each edge shared by at most two
faces, each vertex has a single cyclic link); orientation coherence
(polygon faces have consistent orientation; the signed area is positive for
a counter-clockwise boundary, by the determinant of
Section~\ref{sec:math_orientation_determinants}); and no dangling or
self-intersecting boundaries in polygon results.

A failing invariant check is a bug report with a location; an algorithm that
silently violates an invariant on the next input is a crash waiting for the
input that matters.
\subsection{Common Pitfalls and Bugs to Test For}
\label{sec:appimpl_pitfalls}
A large fraction of defects in geometric programs come from a short list,
and each has a characteristic test:

\begin{itemize}
  \item \index{floating-point comparison}\textbf{Comparing floats with
        \texttt{==}.} Equality on computed coordinates is false in practice.
        Test that snapped or identical inputs produce identical keys, and
        that non-identical inputs are never merged.
  \item \textbf{Vertex ordering.} Mixing counter-clockwise and clockwise
        polygons flips every signed-area and inside-outside decision. Test
        each polygon family in both orientations.
  \item \textbf{Boundary cases.} A point on an edge, a query at a vertex, a
        segment touching another at its endpoint: define half-open rules
        once (e.g., the ray-casting convention of
        Section~\ref{sec:point_in_polygon_winding}) and test every boundary
        case against them.
  \item \textbf{Infinite loops.} Ray-casting that counts a vertex twice, or
        a visibility walk that cycles, hangs forever. Test queries through
        vertices and collinear edges, and assert every walk terminates in a
        bounded number of steps (Section~\ref{sec:point_location_walking}).
  \item \textbf{Sweep-line ordering.} The status of a sweep line
        (Section~\ref{sec:segint_sweep_line}) is ordered by a comparator at
        the current position; when segments share that position the
        comparator must be total and stable or the event queue corrupts.
        Test sweeps with vertical segments, coincident endpoints, and
        segments meeting at endpoints (Section~\ref{sec:segint_degenerate}).
  \item \textbf{Order-of-input dependence.} A permutation of the input must
        not change the output \emph{topology}. Test by shuffling inputs and
        diffing the combinatorial output.
\end{itemize}
\section{Generating Adversarial Test Cases}
\label{sec:app_impl_adversarial}
Degenerate and near-degenerate inputs are where robustness is won or lost.
A dedicated generator that produces them by construction beats hoping random
sampling stumbles on them.
\subsection{Degenerate Configuration Factories}
\label{sec:appimpl_degenerate_factories}
Write small factories that emit the configurations every algorithm fears:
\textbf{coincident points} (duplicates, and duplicates plus noise);
\textbf{collinear sets} (points on a line, vertical and horizontal lines, a
whole dataset collapsed to one line); \textbf{cocircular sets} (points on a
circle, cocircular with a test point exactly on the circumcircle, and
cocircular sets fed into a Delaunay edge flip); \textbf{grids and lattices}
(integer-coordinate inputs that require exactness, including 31-bit
coordinates that push the 64-bit integer bound of
Section~\ref{sec:appimpl_integer}); and \textbf{wide dynamic range}
(coordinates spanning $10^{-9}$ to $10^9$ in one dataset, which break
absolute tolerances and exercise the filter bounds).
\subsection{Near-Degenerate Construction}
\label{sec:appimpl_near_degenerate}
The most insidious inputs are nearly, but not exactly, degenerate. Build
them deliberately: take an exactly degenerate configuration and perturb a
single coordinate by one unit in the last place; rotate a dataset by tiny
angles and verify the output topology is invariant under rigid motion (the
rotation example of Section~\ref{sec:robustness_failures} is the canonical
case~\cite{kettner2008}); scale a degenerate configuration by powers of two
to move it across exponent boundaries. These inputs distinguish a filtered
exact kernel from a threshold-based one.
\subsection{Predicate Fuzzing}
\label{sec:appimpl_fuzzing}
\index{fuzzing}Predicates admit a direct fuzz test that needs no oracle:
generate random point quadruples, evaluate the sign with your fast predicate
and with a reference built on arbitrary-precision integers, and compare. Any
disagreement is a filter bug. Vary the distribution to stress the filter:
uniform points, points very close to a line, large and small coordinates
mixed, and coordinates that are adversarially chosen products and
differences. A predicate that agrees with the reference on $10^6$ cases is
trustworthy; one that has never seen a near-degenerate case is not.
\section{Benchmarking}
\label{sec:app_impl_benchmarking}
Performance work follows a strict order: measure before changing anything,
change what the profiler points at, measure again. Optimizing a structure
that is not the bottleneck makes programs both slow and complicated.
\subsection{Profiling First}
\label{sec:appimpl_profiling}
\index{profiling}Run a sampling profiler on representative inputs before
touching any code. The profile is usually obvious: some fraction of time in
the predicate entry points, the rest in memory traffic around the mesh or
spatial index. If predicates dominate, look at the filter's
\index{floating-point filter!hit rate}hit rate (fraction of calls resolved
without the exact path) and at inlining; if memory traffic dominates,
change the layout (Section~\ref{sec:appimpl_layout}), not the algorithm. A
filter with a 99\% hit rate that costs 4x per call is usually still right;
a filter with a 50\% hit rate is the bottleneck.
\subsection{Benchmarking Methodology}
\label{sec:appimpl_method}
Benchmarks are only as good as their protocol: use fixed, seeded datasets
saved to disk so runs compare across versions; warm up caches, repeat, and
report the median or minimum, not the first run; vary $n$ across orders of
magnitude ($10^3$ to $10^7$) and report per-element cost so asymptotics are
visible; distinguish CPU-bound from memory-bound by scaling (a memory-bound
workload plateaus once data exceeds the cache); and compare against a
known-good reference implementation when one exists
(Section~\ref{sec:app_soft_cgal}) --- being 2x slower than a mature library
is a finding, and so is being 2x faster with a different memory profile.
\subsection{When to Optimize}
\label{sec:appimpl_when}
The right order is: correctness and robustness first, then measurement, then
optimization of the measured bottleneck, then more measurement. Do not tune
a predicate before its sign is certified correct; do not restructure a mesh
layout before a profile shows memory traffic; do not introduce SoA, 64-bit
indices, or parallelism on faith --- each is a complexity tax that should be
paid only when a benchmark shows it buys back more than it costs.
\section{Visualization and Debugging}
\label{sec:app_impl_visualization}
Geometric bugs are spatial, and the fastest way to find most of them is to
look at the geometry. Visualization is a debugging tool as much as a
presentation tool.
\subsection{Visual Debugging}
\label{sec:appimpl_visual}
Build, once, a tiny rendering harness that draws points, segments, polygons,
and mesh edges with colors per category: input points, output vertices,
predicate-sign regions, edges flipped in this step. Two patterns are
especially effective. First, \textbf{color by predicate sign}: plot the
regions of the plane separated by a line and shade nearby points by
$\operatorname{orient2d}$'s sign, so a misclassified near-collinear point is
immediately visible. Second, \textbf{step-through rendering}: for
incremental algorithms (Section~\ref{sec:randomized_incremental}) render
each insertion or flip as a frame, and the first frame where the structure
becomes inconsistent isolates the offending decision.
\subsection{Minimal Reproducers}
\label{sec:appimpl_reproducers}
When a failure appears on a large input, shrink it: remove points and
re-run until the failure disappears; the smallest input that still fails is
the reproducer, and it is usually far more informative than the original
dataset. Automate the search with a delta-debugging loop that splits the
input in half, keeps whichever half still fails, and stops when neither half
does. A reproducer of size three that crashes a Delaunay triangulation is a
bug report; a reproducer of size ten thousand is a mystery. Many robustness
failures are already documented, and matching a reproducer against the
example families of Section~\ref{sec:robustness_failures} often identifies
the root cause immediately~\cite{kettner2008}.
\subsection{Deterministic Builds and Regression Output}
\label{sec:appimpl_regression}
Make the library deterministic: fixed seed for randomized algorithms,
explicit tie-break rules, and no pointer or hash order leaking into output
order. Then compare outputs byte-for-byte across runs and against golden
files for the degenerate test sets of
Section~\ref{sec:appimpl_degenerate_factories}. A golden file that changes
without a corresponding code change is either a bug or an intentionally
documented behavior change --- either way it deserves review, not silence.
Render each golden case to an SVG or image alongside its textual output so
a reviewer can see at a glance what changed and whether the new output is
correct. With deterministic builds, a large randomized suite runs cheaply in
continuous integration, and the combination of property tests, oracle tests,
degenerate factories, and golden files is what keeps robust geometric
libraries robust over years of evolution.
